% ==============================================================================
\section{Методы вычисления QR-разложения на основе ортогональных преобразований}
\label{sec:qr_orthogonal}
% ==============================================================================

Классический процесс ортогонализации Грама–Шмидта можно интерпретировать в матричном виде как последовательное умножение исходной матрицы $\bm{A} =[\bm{a}_1, \dots, \bm{a}_n]$ на верхнетреугольные матрицы справа для получения матрицы с ортогональными столбцами $\bm{Q}$. Таким образом, $\bm{A} \bm{R}_1 \bm{R}_2 \dots \bm{R}_n = \bm{Q}$, откуда $\bm{A} = \bm{Q}\bm{R}$.

Альтернативный подход заключается в применении унитарных (или ортогональных в вещественном случае) преобразований к матрице $\bm{A}$ слева, чтобы привести её к верхнетреугольному виду $\bm{R}$. Это обеспечивает высокую вычислительную устойчивость, так как умножение на унитарные матрицы сохраняет евклидовы нормы и не приводит к накоплению вычислительной ошибки. Требуется найти набор унитарных матриц $\bm{V}_1, \dots, \bm{V}_n$ таких, что
\begin{equation}
    \bm{V}_n \dots \bm{V}_1 \bm{A} = \bm{R}.
    \label{eq:left_orthogonal}
\end{equation}
В данном разделе рассматриваются два класса таких преобразований: матрицы отражения Хаусхолдера и матрицы вращения Гивенса.

% ==============================================================================
\subsection{Отражения Хаусхолдера}
\label{sec:householder}
% ==============================================================================

\begin{tcbdefinition}[Матрица Хаусхолдера]
\label{def:householder}
    Пусть задан вектор $\bm{v} \in \C^n$, такой что $\norm{\bm{v}}_2 = 1$. \textit{Матрицей отражения Хаусхолдера} называется матрица вида
    \begin{equation}
        \bm{H}(\bm{v}) = \bm{I} - 2\bm{v}\bm{v}\H.
        \label{eq:householder_def}
    \end{equation}
\end{tcbdefinition}

\begin{tcblemma}[Свойства матрицы Хаусхолдера]
\label{lem:householder_prop}
    Для любой матрицы Хаусхолдера $\bm{H}(\bm{v})$ справедливы следующие свойства:
    \begin{itemize}
        \item Эрмитовость: $\bm{H}(\bm{v})\H = \bm{H}(\bm{v})$.
        \item Унитарность: $\bm{H}(\bm{v})\H \bm{H}(\bm{v}) = \bm{I}$. Из эрмитовости также следует инволютивность ($\bm{H}(\bm{v})^2 = \bm{I}$).
    \end{itemize}
\end{tcblemma}

Геометрически умножение на матрицу $\bm{H}(\bm{v})$ соответствует отражению вектора относительно гиперплоскости, нормальным вектором к которой является $\bm{v}$. Основная цель применения матриц Хаусхолдера — обнуление заданных компонент вектора.

\begin{tcbtheorem}[Отображение заданного вектора]
\label{thm:householder_map}
    Для любого вектора $\bm{x} \in \C^n$ существуют вектор $\bm{v} \in \C^n$ ($\norm{\bm{v}}_2 = 1$) и скаляр $\gamma \in \C$ ($\abs{\gamma} = 1$) такие, что
    \begin{equation}
        \bm{H}(\bm{v})\bm{x} = \gamma \norm{\bm{x}}_2 \bm{e}_1,
        \label{eq:householder_map}
    \end{equation}
    где $\bm{e}_1 = (1, 0, \dots, 0)\T$ — первый базисный вектор стандартного базиса.
\end{tcbtheorem}

\begin{tcbproof}[теоремы~\ref{thm:householder_map}]
    Пусть $\bm{y} = \gamma \norm{\bm{x}}_2 \bm{e}_1$, причем $\abs{\gamma}=1$, следовательно, $\norm{\bm{x}}_2 = \norm{\bm{y}}_2$. По определению матрицы Хаусхолдера необходимо выполнение равенства:
    \[
        \bm{x} - 2\bm{v}\inner{\bm{v}, \bm{x}} = \bm{y} \implies 2\bm{v}\inner{\bm{v}, \bm{x}} = \bm{x} - \bm{y}.
    \]
    Отсюда следует, что вектор $\bm{v}$ коллинеарен разности $\bm{x} - \bm{y}$. Нормируя его, получаем представление:
    \begin{equation}
        \bm{v} = \frac{\bm{x} - \bm{y}}{\norm{\bm{x} - \bm{y}}_2} = \frac{\bm{x} - \gamma \norm{\bm{x}}_2 \bm{e}_1}{\norm{\bm{x} - \gamma \norm{\bm{x}}_2 \bm{e}_1}_2}.
        \label{eq:v_construction}
    \end{equation}
    Для корректности равенства $\bm{H}(\bm{v})\bm{x} = \bm{y}$ требуется, чтобы скалярное произведение в знаменателе согласовывалось с условиями. Раскрывая квадрат нормы в знаменателе, получаем:
    \[
        \norm{\bm{x} - \bm{y}}_2^2 = \norm{\bm{x}}_2^2 + \norm{\bm{y}}_2^2 - 2\Re\inner{\bm{y}, \bm{x}} = 2\norm{\bm{x}}_2^2 - 2\Re(\bar{\gamma}\norm{\bm{x}}_2 \bar{\bm{e}}_1\H \bm{x}).
    \]
    Замечая, что $\bar{\bm{e}}_1\H \bm{x} = x_1$, условие корректности сводится к тому, что величина $\bar{\gamma} x_1$ должна быть вещественной. Это условие выполняется при выборе $\gamma = \pm \frac{x_1}{\abs{x_1}}$.
\end{tcbproof}

\begin{tcbwarning}[Вычислительная устойчивость (Catastrophic Cancellation)]
    Если выбрать знак $\gamma$ так, что $\gamma \norm{\bm{x}}_2$ будет близко к $x_1$, то при вычислении первой компоненты вектора $\bm{v}$ возникнет вычитание близких чисел ($x_1 - \gamma \norm{\bm{x}}_2 \approx 0$), что приведет к значительной потере точности. Для избежания этой проблемы знак $\gamma$ всегда выбирается противоположным знаку $x_1$:
    \begin{equation}
        \gamma = -\frac{x_1}{\abs{x_1}}.
        \label{eq:gamma_choice}
    \end{equation}
\end{tcbwarning}

\begin{tcbalgorithmbox}[Построение отражения Хаусхолдера]
\label{alg:householder_vec}
    \textbf{Вход:} вектор $\bm{x} \in \C^n$.\\
    \textbf{Выход:} вектор $\bm{v} \in \C^n$ и константа $\gamma$.
    \begin{enumerate}
        \item Вычислить норму $\rho \gets \norm{\bm{x}}_2$.
        \item Вычислить скаляр $\gamma \gets -x_1 / \abs{x_1}$ (если $x_1 = 0$, положить $\gamma = 1$).
        \item Сформировать целевой вектор $\bm{u} \gets \bm{x} - \gamma \rho \bm{e}_1$.
        \item Нормировать вектор $\bm{v} \gets \bm{u} / \norm{\bm{u}}_2$.
        \item Вернуть пару $(\bm{v}, \gamma)$.
    \end{enumerate}
\end{tcbalgorithmbox}

% ==============================================================================
\subsection{Существование QR-разложения}
\label{sec:qr_existence}
% ==============================================================================

\begin{tcbtheorem}[Существование QR-разложения]
\label{thm:qr_exist}
    Для любой матрицы $\bm{A} \in \C^{m \times n}$ существует унитарная матрица $\bm{Q} \in \C^{m \times m}$ и верхнетреугольная матрица $\bm{R} \in \C^{m \times n}$ такие, что $\bm{A} = \bm{Q}\bm{R}$.
\end{tcbtheorem}

\begin{tcbproof}[теоремы~\ref{thm:qr_exist}]
    Доказательство является конструктивным. Применим \cref{alg:householder_vec} к первому столбцу $\bm{a}_1$ матрицы $\bm{A}$ и построим матрицу $\bm{H}_1 = \bm{H}(\bm{v}_1)$. Умножение даст:
    \[
        \bm{H}_1 \bm{A} = 
        \begin{pNiceMatrix}[margin]
            * & * & * & \Cdots & * \\
            0 & * & * & \Cdots & * \\
            \Vdots & \Vdots & \Vdots & & \Vdots \\
            0 & * & * & \Cdots & *
        \end{pNiceMatrix}.
    \]
    Далее рассматривается подматрица без первой строки и первого столбца. Для её первого столбца строится вектор $\bm{v}_2$, и матрица $\bm{H}_2$ формируется в блочном виде:
    \begin{equation}
        \bm{H}_2 = 
        \begin{pNiceMatrix}[margin]
            1 & \bm{0} \\
            \bm{0} & \bm{H}(\bm{v}_2)
        \end{pNiceMatrix}.
        \label{eq:h2_block}
    \end{equation}
    Умножение $\bm{H}_2 \bm{H}_1 \bm{A}$ обнулит поддиагональные элементы второго столбца, не нарушив нули в первом. Продолжая процесс $\min(m-1, n)$ раз, получаем:
    \[
        \bm{H}_k \dots \bm{H}_2 \bm{H}_1 \bm{A} = \bm{R}.
    \]
    Полагая $\bm{Q} = (\bm{H}_k \dots \bm{H}_1)\H = \bm{H}_1 \bm{H}_2 \dots \bm{H}_k$ (в силу эрмитовости матриц Хаусхолдера), получаем требуемое разложение $\bm{A} = \bm{Q}\bm{R}$.
\end{tcbproof}

\begin{tcbremark}[Вычислительные аспекты и «Тонкое» QR-разложение]
    \label{rem:thin_qr}
    \begin{itemize}
        \item Если $m \geq n$, матрица $\bm{R}$ содержит $m-n$ нулевых строк внизу. Блочное перемножение показывает, что $\bm{A} = \bm{Q}_1 \bm{R}_1$, где $\bm{Q}_1$ состоит из первых $n$ столбцов $\bm{Q}$, а $\bm{R}_1$ имеет размер $n \times n$. Это представление называется \textit{тонким QR-разложением} (Thin QR) и используется по умолчанию в большинстве математических пакетов (например, NumPy).
        \item Матрица $\bm{Q}$ никогда не формируется в явном виде путем прямого умножения матриц $\bm{H}_i$. Вместо этого хранится набор векторов $\bm{v}_i$, и умножение $\bm{Q}\bm{x}$ реализуется через последовательное вычитание $\bm{x} \gets \bm{x} - 2\bm{v}_i\inner{\bm{v}_i, \bm{x}}$. Вычислительная сложность получения матрицы $\bm{R}$ составляет $\sim 2mn^2 - \frac{2}{3}n^3$ операций.
    \end{itemize}
\end{tcbremark}

% ==============================================================================
\subsection{Вращения Гивенса}
\label{sec:givens}
% ==============================================================================

Второй подход к ортогональному обнулению элементов опирается на плоские вращения.

\begin{tcbdefinition}[Матрица вращения Гивенса]
\label{def:givens}
    Матрицей вращения Гивенса $\bm{G}_{ij}(\phi) \in \R^{m \times m}$ называется матрица, элементы которой совпадают с элементами единичной матрицы, за исключением подматрицы на пересечении $i$-х и $j$-х строк и столбцов, которая имеет вид:
    \begin{equation}
        \begin{pNiceMatrix}
            \cos\phi & \sin\phi \\
            -\sin\phi & \cos\phi
        \end{pNiceMatrix}.
        \label{eq:givens_block}
    \end{equation}
\end{tcbdefinition}

Умножение вектора на $\bm{G}_{ij}(\phi)$ изменяет только его $i$-ю и $j$-ю компоненты. Подбирая угол $\phi$, можно прицельно обнулить один из этих элементов (например, $j$-й элемент). 

\begin{tcbinsight}[Сравнение Хаусхолдера и Гивенса]
    Асимптотическая сложность QR-разложения через вращения Гивенса составляет $\sim 3mn^2 - n^3$, что формально превышает сложность метода Хаусхолдера. Однако вращения Гивенса предпочтительны в двух сценариях:
    \begin{enumerate}
        \item \textbf{Разреженные матрицы:} Гивенс позволяет обнулять конкретные элементы, сохраняя структуру нулей матрицы, в то время как Хаусхолдер приводит к ее заполнению.
        \item \textbf{Параллельные вычисления:} Преобразование действует локально на пары строк. Непересекающиеся пары строк могут обрабатываться в параллельных потоках независимо друг от друга.
    \end{enumerate}
\end{tcbinsight}

% ==============================================================================
\subsection{QR-разложение с выявлением ранга (Rank-Revealing QR)}
\label{sec:rrqr}
% ==============================================================================

Стандартное QR-разложение последовательно ортогонализует столбцы матрицы слева направо. Если матрица близка к вырожденной (имеет малоранговое приближение), элементы на диагонали $\bm{R}$ могут вести себя непредсказуемо, что не позволяет явно судить о ранге.

Для выявления ранга вводится матрица перестановок столбцов $\bm{P}$. Цель RRQR (Rank-Revealing QR) — найти такое разложение $\bm{A}\bm{P} = \bm{Q}\bm{R}$, при котором блок малых значений скапливается в правом нижнем углу матрицы $\bm{R}$:
\begin{equation}
    \bm{A}\bm{P} = \bm{Q} 
    \begin{pNiceMatrix}[margin]
        \bm{R}_{11} & \bm{R}_{12} \\
        \bm{0}      & \bm{R}_{22}
    \end{pNiceMatrix},
    \label{eq:rrqr_structure}
\end{equation}
где норма блока $\bm{R}_{22}$ минимизируется. Если $\norm{\bm{R}_{22}}$ мала, её можно отбросить, получив скелетное малоранговое разложение $\bm{A} \approx \bm{Q}_1 [\bm{R}_{11}\ \bm{R}_{12}] \bm{P}\T$. 

Задача поиска оптимальной перестановки является NP-трудной, поэтому на практике применяются жадные эвристические алгоритмы.

\begin{tcbalgorithmbox}[Жадный алгоритм RRQR]
\label{alg:rrqr}
    \textbf{Вход:} Матрица $\bm{A} \in \C^{m \times n}$.\\
    \textbf{Выход:} $\bm{Q}, \bm{R}$, матрица перестановок $\bm{P}$.
    \begin{enumerate}
        \item На шаге $k$, пусть активная непереработанная подматрица есть $\bm{R}_{22}^{(k)}$.
        \item Найти столбец в $\bm{R}_{22}^{(k)}$, обладающий наибольшей евклидовой нормой.
        \item Переставить найденный столбец на первое место в $\bm{R}_{22}^{(k)}$ (зафиксировать перестановку в $\bm{P}$).
        \item Применить преобразование Хаусхолдера для обнуления поддиагональных элементов этого столбца. Наибольшая норма концентрируется на диагонали.
        \item Перейти к шагу $k+1$.
    \end{enumerate}
\end{tcbalgorithmbox}

